
%% bare_conf.tex
%% V1.3
%% 2007/01/11
%% by Michael Shell
%% See:
%% http://www.michaelshell.org/
%% for current contact information.
%%
%% This is a skeleton file demonstrating the use of IEEEtran.cls
%% (requires IEEEtran.cls version 1.7 or later) with an IEEE conference paper.
%%
%% Support sites:
%% http://www.michaelshell.org/tex/ieeetran/
%% http://www.ctan.org/tex-archive/macros/latex/contrib/IEEEtran/
%% and
%% http://www.ieee.org/

%%*************************************************************************
%% Legal Notice:
%% This code is offered as-is without any warranty either expressed or
%% implied; without even the implied warranty of MERCHANTABILITY or
%% FITNESS FOR A PARTICULAR PURPOSE!
%% User assumes all risk.
%% In no event shall IEEE or any contributor to this code be liable for
%% any damages or losses, including, but not limited to, incidental,
%% consequential, or any other damages, resulting from the use or misuse
%% of any information contained here.
%%
%% All comments are the opinions of their respective authors and are not
%% necessarily endorsed by the IEEE.
%%
%% This work is distributed under the LaTeX Project Public License (LPPL)
%% ( http://www.latex-project.org/ ) version 1.3, and may be freely used,
%% distributed and modified. A copy of the LPPL, version 1.3, is included
%% in the base LaTeX documentation of all distributions of LaTeX released
%% 2003/12/01 or later.
%% Retain all contribution notices and credits.
%% ** Modified files should be clearly indicated as such, including  **
%% ** renaming them and changing author support contact information. **
%%
%% File list of work: IEEEtran.cls, IEEEtran_HOWTO.pdf, bare_adv.tex,
%%                    bare_conf.tex, bare_jrnl.tex, bare_jrnl_compsoc.tex
%%*************************************************************************

% *** Authors should verify (and, if needed, correct) their LaTeX system  ***
% *** with the testflow diagnostic prior to trusting their LaTeX platform ***
% *** with production work. IEEE's font choices can trigger bugs that do  ***
% *** not appear when using other class files.                            ***
% The testflow support page is at:
% http://www.michaelshell.org/tex/testflow/



% Note that the a4paper option is mainly intended so that authors in
% countries using A4 can easily print to A4 and see how their papers will
% look in print - the typesetting of the document will not typically be
% affected with changes in paper size (but the bottom and side margins will).
% Use the testflow package mentioned above to verify correct handling of
% both paper sizes by the user's LaTeX system.
%
% Also note that the "draftcls" or "draftclsnofoot", not "draft", option
% should be used if it is desired that the figures are to be displayed in
% draft mode.
%
\documentclass[10pt, conference, a4paper, compsocconf]{IEEEtran}
% customized to Vietnamese codes and free-modifying header/footer by Nguyễn Trường Thắng
\usepackage{amsmath,amsxtra,amssymb,amsthm,latexsym,amscd,amsfonts}
\usepackage[utf8]{vietnam}
%% We will completely define our own header strings, so switch the fancy
%% headers on, but nuke all the default values. (Note that this package
%% *has* to load after the geometry package!)


%%% Edited 22/3/2015 (NHA)

\usepackage[top=2.7cm, bottom = 5.4cm, left=2.72cm, right = 3.0cm]{geometry}
\headsep = 0.7cm
\headheight = 1.0cm 
\footskip = 1.0cm
\voffset = -0.74 cm 
%%%%

\usepackage{fancyhdr}
\pagestyle{fancy}
%\fancyhf{}



%% Now begin customising things. See the fancyhdr docs for more info.

%\renewcommand{\chaptermark}[1]{\markboth{\MakeUppercase{#1}}{}}
\renewcommand{\sectionmark}[1]{\markright{\MakeUppercase{#1}}{}}
\renewcommand{\headrulewidth}{0pt}
% end of customization

% Add the compsocconf option for Computer Society conferences.
%
% If IEEEtran.cls has not been installed into the LaTeX system files,
% manually specify the path to it like:
% \documentclass[conference]{../sty/IEEEtran}


% Some very useful LaTeX packages include:
% (uncomment the ones you want to load)


% *** MISC UTILITY PACKAGES ***
%
%\usepackage{ifpdf}
% Heiko Oberdiek's ifpdf.sty is very useful if you need conditional
% compilation based on whether the output is pdf or dvi.
% usage:
% \ifpdf
%   % pdf code
% \else
%   % dvi code
% \fi
% The latest version of ifpdf.sty can be obtained from:
% http://www.ctan.org/tex-archive/macros/latex/contrib/oberdiek/
% Also, note that IEEEtran.cls V1.7 and later provides a builtin
% \ifCLASSINFOpdf conditional that works the same way.
% When switching from latex to pdflatex and vice-versa, the compiler may
% have to be run twice to clear warning/error messages.






% *** CITATION PACKAGES ***
%
\usepackage{cite}
% cite.sty was written by Donald Arseneau
% V1.6 and later of IEEEtran pre-defines the format of the cite.sty package
% \cite{} output to follow that of IEEE. Loading the cite package will
% result in citation numbers being automatically sorted and properly
% "compressed/ranged". e.g., [1], [9], [2], [7], [5], [6] without using
% cite.sty will become [1], [2], [5]--[7], [9] using cite.sty. cite.sty's
% \cite will automatically add leading space, if needed. Use cite.sty's
% noadjust option (cite.sty V3.8 and later) if you want to turn this off.
% cite.sty is already installed on most LaTeX systems. Be sure and use
% version 4.0 (2003-05-27) and later if using hyperref.sty. cite.sty does
% not currently provide for hyperlinked citations.
% The latest version can be obtained at:
% http://www.ctan.org/tex-archive/macros/latex/contrib/cite/
% The documentation is contained in the cite.sty file itself.






% *** GRAPHICS RELATED PACKAGES ***
%
\ifCLASSINFOpdf
\usepackage[pdftex]{graphicx}
  % declare the path(s) where your graphic files are
  % \graphicspath{{../pdf/}{../jpeg/}}
  % and their extensions so you won't have to specify these with
  % every instance of \includegraphics
  \DeclareGraphicsExtensions{.pdf,.jpeg,.png,.jpg}
\else
  % or other class option (dvipsone, dvipdf, if not using dvips). graphicx
  % will default to the driver specified in the system graphics.cfg if no
  % driver is specified.
\usepackage[dvips]{graphicx}
  % declare the path(s) where your graphic files are
  % \graphicspath{{../eps/}}
  % and their extensions so you won't have to specify these with
  % every instance of \includegraphics
  \DeclareGraphicsExtensions{.eps}
\fi
% graphicx was written by David Carlisle and Sebastian Rahtz. It is
% required if you want graphics, photos, etc. graphicx.sty is already
% installed on most LaTeX systems. The latest version and documentation can
% be obtained at:
% http://www.ctan.org/tex-archive/macros/latex/required/graphics/
% Another good source of documentation is "Using Imported Graphics in
% LaTeX2e" by Keith Reckdahl which can be found as epslatex.ps or
% epslatex.pdf at: http://www.ctan.org/tex-archive/info/
%
% latex, and pdflatex in dvi mode, support graphics in encapsulated
% postscript (.eps) format. pdflatex in pdf mode supports graphics
% in .pdf, .jpeg, .png and .mps (metapost) formats. Users should ensure
% that all non-photo figures use a vector format (.eps, .pdf, .mps) and
% not a bitmapped formats (.jpeg, .png). IEEE frowns on bitmapped formats
% which can result in "jaggedy"/blurry rendering of lines and letters as
% well as large increases in file sizes.
%
% You can find documentation about the pdfTeX application at:
% http://www.tug.org/applications/pdftex


% *** MATH PACKAGES ***
%
%\usepackage[cmex10]{amsmath}
% A popular package from the American Mathematical Society that provides
% many useful and powerful commands for dealing with mathematics. If using
% it, be sure to load this package with the cmex10 option to ensure that
% only type 1 fonts will utilized at all point sizes. Without this option,
% it is possible that some math symbols, particularly those within
% footnotes, will be rendered in bitmap form which will result in a
% document that can not be IEEE Xplore compliant!
%
% Also, note that the amsmath package sets \interdisplaylinepenalty to 10000
% thus preventing page breaks from occurring within multiline equations. Use:
%\interdisplaylinepenalty=2500
% after loading amsmath to restore such page breaks as IEEEtran.cls normally
% does. amsmath.sty is already installed on most LaTeX systems. The latest
% version and documentation can be obtained at:
% http://www.ctan.org/tex-archive/macros/latex/required/amslatex/math/


% *** SPECIALIZED LIST PACKAGES ***
%
%\usepackage{algorithmic}
% algorithmic.sty was written by Peter Williams and Rogerio Brito.
% This package provides an algorithmic environment fo describing algorithms.
% You can use the algorithmic environment in-text or within a figure
% environment to provide for a floating algorithm. Do NOT use the algorithm
% floating environment provided by algorithm.sty (by the same authors) or
% algorithm2e.sty (by Christophe Fiorio) as IEEE does not use dedicated
% algorithm float types and packages that provide these will not provide
% correct IEEE style captions. The latest version and documentation of
% algorithmic.sty can be obtained at:
% http://www.ctan.org/tex-archive/macros/latex/contrib/algorithms/
% There is also a support site at:
% http://algorithms.berlios.de/index.html
% Also of interest may be the (relatively newer and more customizable)
% algorithmicx.sty package by Szasz Janos:
% http://www.ctan.org/tex-archive/macros/latex/contrib/algorithmicx/


% *** ALIGNMENT PACKAGES ***
%
\usepackage{array}
% Frank Mittelbach's and David Carlisle's array.sty patches and improves
% the standard LaTeX2e array and tabular environments to provide better
% appearance and additional user controls. As the default LaTeX2e table
% generation code is lacking to the point of almost being broken with
% respect to the quality of the end results, all users are strongly
% advised to use an enhanced (at the very least that provided by array.sty)
% set of table tools. array.sty is already installed on most systems. The
% latest version and documentation can be obtained at:
% http://www.ctan.org/tex-archive/macros/latex/required/tools/


%\usepackage{mdwmath}
%\usepackage{mdwtab}
% Also highly recommended is Mark Wooding's extremely powerful MDW tools,
% especially mdwmath.sty and mdwtab.sty which are used to format equations
% and tables, respectively. The MDWtools set is already installed on most
% LaTeX systems. The lastest version and documentation is available at:
% http://www.ctan.org/tex-archive/macros/latex/contrib/mdwtools/


% IEEEtran contains the IEEEeqnarray family of commands that can be used to
% generate multiline equations as well as matrices, tables, etc., of high
% quality.


%\usepackage{eqparbox}
% Also of notable interest is Scott Pakin's eqparbox package for creating
% (automatically sized) equal width boxes - aka "natural width parboxes".
% Available at:
% http://www.ctan.org/tex-archive/macros/latex/contrib/eqparbox/


% *** SUBFIGURE PACKAGES ***
%\usepackage[tight,footnotesize]{subfigure}
% subfigure.sty was written by Steven Douglas Cochran. This package makes it
% easy to put subfigures in your figures. e.g., "Figure 1a and 1b". For IEEE
% work, it is a good idea to load it with the tight package option to reduce
% the amount of white space around the subfigures. subfigure.sty is already
% installed on most LaTeX systems. The latest version and documentation can
% be obtained at:
% http://www.ctan.org/tex-archive/obsolete/macros/latex/contrib/subfigure/
% subfigure.sty has been superceeded by subfig.sty.



%\usepackage[caption=false]{caption}
\usepackage[font=footnotesize]{subfig}
% subfig.sty, also written by Steven Douglas Cochran, is the modern
% replacement for subfigure.sty. However, subfig.sty requires and
% automatically loads Axel Sommerfeldt's caption.sty which will override
% IEEEtran.cls handling of captions and this will result in nonIEEE style
% figure/table captions. To prevent this problem, be sure and preload
% caption.sty with its "caption=false" package option. This is will preserve
% IEEEtran.cls handing of captions. Version 1.3 (2005/06/28) and later
% (recommended due to many improvements over 1.2) of subfig.sty supports
% the caption=false option directly:
%\usepackage[caption=false,font=footnotesize]{subfig}
%
% The latest version and documentation can be obtained at:
% http://www.ctan.org/tex-archive/macros/latex/contrib/subfig/
% The latest version and documentation of caption.sty can be obtained at:
% http://www.ctan.org/tex-archive/macros/latex/contrib/caption/




% *** FLOAT PACKAGES ***
%
%\usepackage{fixltx2e}
% fixltx2e, the successor to the earlier fix2col.sty, was written by
% Frank Mittelbach and David Carlisle. This package corrects a few problems
% in the LaTeX2e kernel, the most notable of which is that in current
% LaTeX2e releases, the ordering of single and double column floats is not
% guaranteed to be preserved. Thus, an unpatched LaTeX2e can allow a
% single column figure to be placed prior to an earlier double column
% figure. The latest version and documentation can be found at:
% http://www.ctan.org/tex-archive/macros/latex/base/



%\usepackage{stfloats}
% stfloats.sty was written by Sigitas Tolusis. This package gives LaTeX2e
% the ability to do double column floats at the bottom of the page as well
% as the top. (e.g., "\begin{figure*}[!b]" is not normally possible in
% LaTeX2e). It also provides a command:
%\fnbelowfloat
% to enable the placement of footnotes below bottom floats (the standard
% LaTeX2e kernel puts them above bottom floats). This is an invasive package
% which rewrites many portions of the LaTeX2e float routines. It may not work
% with other packages that modify the LaTeX2e float routines. The latest
% version and documentation can be obtained at:
% http://www.ctan.org/tex-archive/macros/latex/contrib/sttools/
% Documentation is contained in the stfloats.sty comments as well as in the
% presfull.pdf file. Do not use the stfloats baselinefloat ability as IEEE
% does not allow \baselineskip to stretch. Authors submitting work to the
% IEEE should note that IEEE rarely uses double column equations and
% that authors should try to avoid such use. Do not be tempted to use the
% cuted.sty or midfloat.sty packages (also by Sigitas Tolusis) as IEEE does
% not format its papers in such ways.


% *** PDF, URL AND HYPERLINK PACKAGES ***
%
%\usepackage{url}
% url.sty was written by Donald Arseneau. It provides better support for
% handling and breaking URLs. url.sty is already installed on most LaTeX
% systems. The latest version can be obtained at:
% http://www.ctan.org/tex-archive/macros/latex/contrib/misc/
% Read the url.sty source comments for usage information. Basically,
% \url{my_url_here}.


% *** Do not adjust lengths that control margins, column widths, etc. ***
% *** Do not use packages that alter fonts (such as pslatex).         ***
% There should be no need to do such things with IEEEtran.cls V1.6 and later.
% (Unless specifically asked to do so by the journal or conference you plan
% to submit to, of course. )


% correct bad hyphenation here
%\hyphenation{op-tical net-works semi-conduc-tor}

%%% Edited 6/6/2014
% title page
%\makeatletter
%\def\ps@IEEEtitlepagestyle{%
%\def\@oddhead{\mbox{\centering{\footnotesize{\textit{Hội thảo quốc gia lần thứ XVIII: Một số vấn đề chọn lọc của Công nghệ thông tin và truyền thông-TP HCM, 05-06/11/2015}}\scriptsize\rightmark \hfil \thepage}}%
%\def\@evenhead{\scriptsize\thepage \hfil \leftmark\mbox{\centering{\footnotesize{\textit{Hội thảo quốc gia lần thứ XVII: Một số vấn đề chọn lọc của Công nghệ thông tin và truyền thông Đắk Lắk, 30-31/10/2014}}}}}%
%}
%}
%\makeatother

\makeatletter
\def\ps@IEEEtitlepagestyle{%
\def\@oddhead{\centering{\mbox{\footnotesize{\textit{\;\;\;\; Hội thảo quốc gia lần thứ XXIX: Một số vấn đề chọn lọc của Công nghệ thông tin và truyền thông -- Hà Nội, 7-8/11/2026}}}}%
\def\@evenhead{\centering{\mbox{\footnotesize{\textit{\;\;\;\; Hội thảo quốc gia lần thứ XXIX: Một số vấn đề chọn lọc của Công nghệ thông tin và truyền thông -- Hà Nội, 7-8/11/2026}}}}}%
\def\@oddfoot{\scriptsize \thepage \hfil }%
\def\@evenfoot{\scriptsize \hfil \thepage}
}
}
\makeatother



\begin{document}
\pagenumbering{gobble}
\fontsize{10}{12}
\selectfont

% customizing header/footer via fancyhdr by Nguyễn Trường Thắng
%% Configuration of the header strings for the frontmatter pages.
%\fancyhead[RO]{{\footnotesize\rightmark}ABC\hspace{2em}\thepage}
%\setcounter{tocdepth}{2}
%\fancyhead[LE]{\thepage\hspace{2em}ABC\footnotesize{\leftmark}}
%\fancyhead[RE,LO]{}
%\fancyhead[RO]{{\footnotesize\rightmark}\hspace{2em}\thepage}

%% Configuration of the header strings for the main manuscript pages.
\fancyhead[RE,LO]{\centering{\footnotesize{\textit{Hội thảo quốc gia lần thứ XXIX: Một số vấn đề chọn lọc của Công nghệ thông tin và truyền thông -- Hà Nội, 7-8/11/2026}}}}
% end of customization

%
% paper title
% can use linebreaks \\ within to get better formatting as desired
\title{An application of GenAi in the problem of detecting logical flaws in smart contracts}


% author names and affiliations
% use a multiple column layout for up to two different
% affiliations

%%% Edited 01/01/2022 (NHA) %%%
\author{\IEEEauthorblockN{Nguyễn Văn A\IEEEauthorrefmark{1}, Nguyễn Văn B\IEEEauthorrefmark{2}, Nguyễn Văn C\IEEEauthorrefmark{3}}
	\IEEEauthorblockA{\IEEEauthorrefmark{1}Viện Công nghệ thông tin, Viện Hàn lâm Khoa học và công nghệ Việt Nam,}
	\IEEEauthorblockA{\IEEEauthorrefmark{2}Học viện Khoa học và công nghệ, \IEEEauthorrefmark{3}Trường Đại học Hà Nội}
	\IEEEauthorblockA{Email: nguyen\_van\_a@abc.xyz, nguyen\_van\_b@abc.xyz, nguyen\_van\_c@abc.xyz}
}

% conference papers do not typically use \thanks and this command
% is locked out in conference mode. If really needed, such as for
% the acknowledgment of grants, issue a \IEEEoverridecommandlockouts
% after \documentclass

% for over three affiliations, or if they all won't fit within the width
% of the page, use this alternative format:
%
%\author{\IEEEauthorblockN{Michael Shell\IEEEauthorrefmark{1},
%Homer Simpson\IEEEauthorrefmark{2},
%James Kirk\IEEEauthorrefmark{3},
%Montgomery Scott\IEEEauthorrefmark{3} and
%Eldon Tyrell\IEEEauthorrefmark{4}}
%\IEEEauthorblockA{\IEEEauthorrefmark{1}School of Electrical and Computer Engineering\\
%Georgia Institute of Technology,
%Atlanta, Georgia 30332--0250\\ Email: see http://www.michaelshell.org/contact.html}
%\IEEEauthorblockA{\IEEEauthorrefmark{2}Twentieth Century Fox, Springfield, USA\\
%Email: homer@thesimpsons.com}
%\IEEEauthorblockA{\IEEEauthorrefmark{3}Starfleet Academy, San Francisco, California 96678-2391\\
%Telephone: (800) 555--1212, Fax: (888) 555--1212}
%\IEEEauthorblockA{\IEEEauthorrefmark{4}Tyrell Inc., 123 Replicant Street, Los Angeles, California 90210--4321}}


% use for special paper notices
%\IEEEspecialpapernotice{(Invited Paper)}




% make the title area
\maketitle


\begin{abstract}
Smart contracts are immutable programs deployed on blockchain networks that directly manage financial assets, making security vulnerabilities irreversible and potentially catastrophic. Existing automated audit tools --- primarily static analyzers such as Slither and Mythril --- suffer from high false positive rates and fail to detect semantic and business-logic flaws. Single-LLM approaches such as GPTScan improve detection coverage but rely on a single analytical perspective without exploitability validation. This paper proposes MECAP (Multi-Expert Contract Audit Panel), a novel multi-agent LLM framework that coordinates 22 specialized agents across five security domains to audit Solidity smart contracts from diverse expert perspectives simultaneously. MECAP introduces three key contributions: (1) a Contract Knowledge Graph that grounds agent reasoning in concrete code evidence, reducing hallucination; (2) a structured three-phase audit session enabling systematic cross-domain challenge and validation of findings; and (3) a three-layer weighted consensus mechanism combining intra-domain agreement, cross-domain validation, and attacker corroboration to filter false positives. Evaluated on the SmartBugs Curated benchmark of 143 labeled contracts, MECAP surpasses static analysis baselines in both false positive rate and logic vulnerability coverage, with demonstrated detection of DeFi-specific vulnerabilities missed by existing tools.
\end{abstract}

\begin{IEEEkeywords}
smart contract; vulnerability detection; multi-agent LLM; knowledge graph; consensus mechanism
\end{IEEEkeywords}


% For peer review papers, you can put extra information on the cover
% page as needed:
% \ifCLASSOPTIONpeerreview
% \begin{center} \bfseries EDICS Category: 3-BBND \end{center}
% \fi
%
% For peerreview papers, this IEEEtran command inserts a page break and
% creates the second title. It will be ignored for other modes.
\IEEEpeerreviewmaketitle



\section{Introduction}

Smart contracts are self-executing programs deployed on blockchain networks that autonomously govern financial transactions without intermediaries. Unlike conventional software, deployed smart contracts are immutable --- once on-chain, their code cannot be modified, meaning any security vulnerability is permanent and its consequences irreversible. The 2016 DAO attack exploited a reentrancy flaw to drain over \$50 million USD \cite{mehar2019dao}, and cumulative losses from DeFi exploits have since surpassed \$10 billion USD \cite{carpentier2025defi}. This combination of immutability and direct financial custody makes pre-deployment security auditing a critical necessity.

Existing automated audit tools fall into two categories, each with fundamental limitations. Static analysis tools such as Slither \cite{feist2019slither} and Mythril \cite{mythril} operate through predefined vulnerability patterns, achieving high throughput but generating substantial false positives and missing semantic flaws entirely. An empirical study across 47,587 Ethereum contracts found that existing tools flagged 97\% of contracts as vulnerable, indicating severe false positive rates \cite{durieux2020empirical}. These tools rely on syntactic pattern matching and lack the semantic reasoning needed to detect business-logic vulnerabilities such as price oracle manipulation or governance attacks. The second category --- single-LLM approaches exemplified by GPTScan \cite{sun2024gptscan} --- improves logic vulnerability coverage but inherits a single analytical perspective and provides no exploitability validation, leaving practitioners unable to distinguish theoretical issues from practically exploitable flaws.

We propose \textbf{MECAP} (Multi-Expert Contract Audit Panel), a novel multi-agent LLM framework that addresses these limitations through structured expert debate. MECAP coordinates 22 specialized agents across five security domains --- Application Security, Blockchain Security, Cryptography, DeFi Protocol Security, and Governance --- organized into two tiers: 17 domain expert agents and 5 attacker profile agents. All agents operate within a Contract Knowledge Graph that grounds their reasoning in concrete code evidence. A three-phase audit session (intra-domain analysis, cross-domain validation, and attacker challenge) systematically surfaces and filters findings. A three-layer weighted consensus mechanism then synthesizes agent outputs into a ranked, confidence-scored vulnerability report with patch suggestions.

The contributions of this paper are as follows:
\begin{itemize}
  \item \textbf{C1 --- Novel Framework:} MECAP is, to the best of our knowledge, the first multi-domain, multi-persona agent panel applied to smart contract vulnerability detection, extending the OASIS social simulation architecture \cite{yang2024oasis} to the security audit domain.
  \item \textbf{C2 --- KG Grounding:} A Contract Knowledge Graph anchors all agent findings to specific functions, state variables, and call sequences, providing verifiable evidence and measurably reducing hallucination compared to context-free LLM analysis.
  \item \textbf{C3 --- Exploitability Validation:} Five attacker profiles (Reentrancy Exploiter, Flash Loan Attacker, Governance Attacker, Access Control Exploiter, Logic Exploiter) distinguish practically exploitable vulnerabilities from theoretical ones --- a capability absent in all existing automated tools.
  \item \textbf{C4 --- Consensus-Based FP Reduction:} A three-layer weighted consensus (intra-domain $w{=}0.30$, cross-domain $w{=}0.45$, attacker corroboration $w{=}0.25$, threshold $=0.35$) filters false positives more effectively than single-threshold rule-based approaches.
  \item \textbf{C5 --- DeFi Coverage:} Dedicated DeFi and Governance domain agents detect flash loan, oracle manipulation, and governance attacks that fall outside the detection scope of static analysis tools.
\end{itemize}

\section{Related Work}

\textbf{Static Analysis Tools.}
Slither \cite{feist2019slither} applies data flow and taint analysis to Solidity source code, detecting common vulnerability patterns at high speed. Mythril \cite{mythril} uses symbolic execution and SMT solving for deeper path analysis but incurs significant runtime overhead on complex contracts \cite{ghaleb2020effective}. A large-scale empirical study across 47,587 Ethereum contracts revealed that these tools produce severe false positive rates, with 97\% of contracts flagged as potentially vulnerable \cite{durieux2020empirical}. Both tools rely on syntactic pattern matching and fundamentally lack the semantic reasoning needed to detect business-logic vulnerabilities such as price oracle manipulation or governance attacks.

\textbf{LLM-Based Approaches.}
GPTScan \cite{sun2024gptscan} combines GPT with static AST analysis to match logic vulnerability scenarios against predefined templates, demonstrating strong precision on token contracts. However, it relies on a single LLM call and fixed templates, with no mechanism for exploitability validation. AuditGPT \cite{xia2024auditgpt} applies LLMs to audit smart contracts against ERC compliance rules but adopts the same single-agent, single-pass architecture. SmartInv \cite{wang2024smartinv} uses multimodal LLM to synthesize contract invariants for correctness verification, providing a complementary approach to vulnerability discovery. All LLM-based approaches to date operate from a single analytical perspective without cross-domain debate or attacker-perspective corroboration.

\textbf{Multi-Agent LLM Systems.}
Du et al.\ \cite{du2023debate} demonstrated that multi-agent debate improves factual accuracy over single-agent inference, providing theoretical grounding for adversarial agent architectures in verification tasks. CAMEL \cite{li2023camel} introduced role-playing cooperative agents, while AgentVerse \cite{chen2023agentverse} explored emergent behaviors in structured multi-agent collaboration. OASIS \cite{yang2024oasis} scaled agent-based social simulation to one million agents through structured interaction rounds --- the architecture that MECAP adapts for security auditing. Knowledge Graph grounding via Retrieval-Augmented Generation \cite{lewis2020rag} reduces LLM hallucination by anchoring responses to external structured knowledge, a principle central to MECAP's Contract KG design.

\textbf{Research Gap.}
No existing approach combines multi-domain expert debate, Knowledge Graph grounding, and attacker-perspective exploitability validation in a unified smart contract audit pipeline. MECAP addresses this gap.

\section{Mothodology}
\subsection{Selecting a Template (Heading 2)}
First, confirm that you have the correct template for your paper size. This template has been tailored for output on the US-letter paper size. If you are using A4-sized paper, please close this template and download the file for A4 paper format called \verb+"CPS_A4_format”+.
\subsection{Maintaining the Integrity of the Specifications}
The template is used to format your paper and style the text. All margins, column widths, line spaces, and text fonts are prescribed; please do not alter them. You may note peculiarities. For example, the head margin in this template measures proportionately more than is customary. This measurement and others are deliberate, using specifications that anticipate your paper as one part of the entire proceedings, and not as an independent document. Please do not revise any of the current designations.

\section{Implementation}
Before you begin to format your paper, first write and save the content as a separate text file. Keep your text and graphic files separate until after the text has been formatted and styled. Do not use hard tabs, and limit use of hard returns to only one return at the end of a paragraph. Do not add any kind of pagination anywhere in the paper. Do not number text heads-the template will do that for you.


Finally, complete content and organizational editing before formatting. Please take note of the following items when proofreading spelling and grammar:
\subsection{Abbreviations and Acronyms}
Define abbreviations and acronyms the first time they are used in the text, even after they have been defined in the abstract. Abbreviations such as IEEE, SI, MKS, CGS, sc, dc, and rms do not have to be defined. Do not use abbreviations in the title or heads unless they are unavoidable.
\subsection{Units}
\begin{itemize}
\item Use either SI (MKS) or CGS as primary units. (SI units are encouraged.) English units may be used as secondary units (in parentheses). An exception would be the use of English units as identifiers in trade, such as “3.5-inch disk drive”.
\item Avoid combining SI and CGS units, such as current in amperes and magnetic field in oersteds. This often leads to confusion because equations do not balance dimensionally. If you must use mixed units, clearly state the units for each quantity that you use in an equation.
\item Do not mix complete spellings and abbreviations of units: “Wb/m2” or “webers per square meter”, not “webers/m2”.  Spell out units when they appear in text: “. . . a few henries”, not “. . . a few H”.
\item Use a zero before decimal points: “0.25”, not “.25”.
\end{itemize}
\subsection{Equations}
The equations are an exception to the prescribed specifications of this template. You will need to determine whether or not your equation should be typed using either the Times New Roman or the Symbol font (please no other font). To create multileveled equations, it may be necessary to treat the equation as a graphic and insert it into the text after your paper is styled.

Number equations consecutively. Equation numbers, within parentheses, are to position flush right, as in (1), using a right tab stop. To make your equations more compact, you may use the solidus ( / ), the exp function, or appropriate exponents. Italicize Roman symbols for quantities and variables, but not Greek symbols. Use a long dash rather than a hyphen for a minus sign. Punctuate equations with commas or periods when they are part of a sentence, as in

Note that the equation is centered using a center tab stop. Be sure that the symbols in your equation have been defined before or immediately following the equation. Use “(1)”, not “Eq. (1)” or “equation (1)”, except at the beginning of a sentence: “Equation (1) is . . .”
\subsection{Some Common Mistakes}
\begin{itemize}
\item The word “data” is plural, not singular.
\item The subscript for the permeability of vacuum 0, and other common scientific constants, is zero with subscript formatting, not a lowercase letter “o”.
\item In American English, commas, semi-/colons, periods, question and exclamation marks are located within quotation marks only when a complete thought or name is cited, such as a title or full quotation. When quotation marks are used, instead of a bold or italic typeface, to highlight a word or phrase, punctuation should appear outside of the quotation marks. A parenthetical phrase or statement at the end of a sentence is punctuated outside of the closing parenthesis (like this). (A parenthetical sentence is punctuated within the parentheses.)
\item A graph within a graph is an “inset”, not an “insert”. The word alternatively is preferred to the word “alternately” (unless you really mean something that alternates).
\item Do not use the word “essentially” to mean “approximately” or “effectively”.
\item In your paper title, if the words “that uses” can accurately replace the word “using”, capitalize the “u”; if not, keep using lower-cased.
\item Be aware of the different meanings of the homophones “affect” and “effect”, “complement” and “compliment”, “discreet” and “discrete”, “principal” and “principle”.
\item Do not confuse “imply” and “infer”.
\item The prefix “non” is not a word; it should be joined to the word it modifies, usually without a hyphen.
\item There is no period after the “et” in the Latin abbreviation “et al.”.
\item The abbreviation “i.e.” means “that is”, and the abbreviation “e.g.” means “for example”.
\end{itemize}

\section{Conclusion}
The conclusion goes here. this is more of the conclusion

% conference papers do not normally have an appendix


% use section* for acknowledgement
\section*{Cảm tạ}


The authors would like to thank...
more thanks here


% trigger a \newpage just before the given reference
% number - used to balance the columns on the last page
% adjust value as needed - may need to be readjusted if
% the document is modified later
%\IEEEtriggeratref{8}
% The "triggered" command can be changed if desired:
%\IEEEtriggercmd{\enlargethispage{-5in}}

% references section

% can use a bibliography generated by BibTeX as a .bbl file
% BibTeX documentation can be easily obtained at:
% http://www.ctan.org/tex-archive/biblio/bibtex/contrib/doc/
% The IEEEtran BibTeX style support page is at:
% http://www.michaelshell.org/tex/ieeetran/bibtex/
\bibliographystyle{IEEEtran}
\bibliography{references}




% that's all folks
\end{document}


